% Fuente: Auxiliar 2, P3 (pauta).
\textbf{Parámetros y conjuntos}

\begin{itemize}
    \item $r_i$: ranking del postulante $i$.
    \item $N$: número total de becas.
    \item $n_s$: número de postulantes por estado $s \in \{1,\ldots,50\}$.
    \item $T(s)$: objetivo por estado $s \in \{1,\ldots,50\}$.
    \item $F_s$: conjunto de postulantes del estado $s$.
    \item $T_g$: objetivo de género $g \in \{M,F\}$.
    \item $F_g$: conjunto de postulantes de género $g$.
    \item $T(d)$: objetivo por disciplina $d \in \{1,\ldots,19\}$.
    \item $F_d$: conjunto de postulantes a la disciplina $d$.
    \item $T(l)$: objetivo por nivel de grado $l \in \{u,g1,g2\}$.
    \item $F_l$: conjunto de postulantes de nivel de grado $l$.
\end{itemize}

\textbf{Variables de decisión}

\begin{equation}
x_i =
\begin{cases}
1, & \text{si el postulante } i \text{ recibe la beca} \\
0, & \text{en otro caso}
\end{cases}
\end{equation}

\textbf{Restricciones}

\begin{itemize}
    \item Como todos los postulantes del grupo 1 ganaron la beca,
    \begin{equation}
    x_i = 1, \quad \forall i \in G1
    \end{equation}

    \item La suma de las variables de decisión debe ser igual al total de becas:
    \begin{equation}
    \sum_{i=1}^{|G1| + K} x_i = N
    \end{equation}

    \item Se debe cumplir el objetivo de género:
    \begin{equation}
    \sum_{i \in F_g} x_i \geq T(g), \quad \forall g \in \{M,F\}
    \end{equation}

    \item Se debe cumplir el objetivo de disciplina:
    \begin{equation}
    \sum_{i \in F_d} x_i \geq T(d), \quad \forall d \in \{1,\ldots,19\}
    \end{equation}

    \item Se debe cumplir el objetivo de nivel de grado:
    \begin{equation}
    \sum_{i \in F_l} x_i \geq T(l), \quad \forall l \in \{u,g1,g2\}
    \end{equation}

    \item Se debe cumplir el objetivo de estado, cuidando que no se supere el número de postulantes por estado:
    \begin{equation}
    \sum_{i \in F_s} x_i \geq \min(T(s),n_s), \quad \forall s \in \{1,\ldots,50\}
    \end{equation}
\end{itemize}

\textbf{Función objetivo}

Se requiere minimizar el mayor (peor) ranking entre los postulantes de las becas, entonces
la función objetivo es:

\begin{equation}
\min \max_{i \in G2} r_i x_i
\end{equation}

Esta última expresión se puede reformular como:

\begin{equation}
\min z
\end{equation}

\begin{equation}
z \geq r_i x_i, \quad \forall i \in G2
\end{equation}

\textbf{Modelo completo}

\begin{equation}
\min z
\end{equation}

\begin{equation}
z \geq r_i x_i, \quad \forall i \in G2
\end{equation}

\begin{equation}
x_i = 1, \quad \forall i \in G1
\end{equation}

\begin{equation}
\sum_{i=1}^{|G1| + K} x_i = N
\end{equation}

\begin{equation}
\sum_{i \in F_g} x_i \geq T(g), \quad \forall g \in \{M,F\}
\end{equation}

\begin{equation}
\sum_{i \in F_d} x_i \geq T(d), \quad \forall d \in \{1,\ldots,19\}
\end{equation}

\begin{equation}
\sum_{i \in F_l} x_i \geq T(l), \quad \forall l \in \{u,g1,g2\}
\end{equation}

\begin{equation}
\sum_{i \in F_s} x_i \geq \min(T(s),n_s), \quad \forall s \in \{1,\ldots,50\}
\end{equation}

\begin{equation}
x_i \in \{0,1\} \quad \forall i \in G1
\end{equation}
